% Adam report draft; prepared with AI assistance for author review.
Our results show that field of study is closely related to reported grades, and that these differences can affect access when students are ranked together. Shared field effects explain 65\% of the variation in Spanish field means within university-years; this is not a percentage of individual students' grades. In our simulation, the median share reaching the top 20\% is 41\% for Nursing and 6\% for Business administration. In Portugal, within-degree percentiles reduce the spread of selection rates from 0.090 to 0.006, with a small difference in later grades: 14.21 versus 14.08. The instructor's adjusted mark reduces the graduate-weighted distance from equal access from 11.8 to 5.1 percentage points, although differences in spread remain. Using universities' past field profiles also lowers test prediction error from 0.225 to 0.138 grade points. The US results suggest a similar broad-area pattern, but remain uncertain.

For institutions aiming to give students comparable chances across degrees, within-degree percentiles are worth considering. They account for differences in both average grades and their spread, while subtracting a field deviation changes only the average. ECTS grade-distribution tables offer a related way to put grades in context, although they do not prescribe a selection rule (European Commission 39--41). Institutions still need to decide what their ranking should reward. Our results describe associations between fields and grades; they do not establish that grading bias causes the gaps.

The next step is to test these ideas closer to actual decisions. Records of students taking shared courses, ideally with common assessments, could help separate preparation from grading differences. Individual Spanish records would let us examine real rankings rather than simulated selection shares. Finally, an exchange-allocation pilot could compare selection rules using later academic performance and participation across degrees, with the criteria agreed beforehand. This would help assess both the access benefits and possible costs that one semester's grades cannot capture.
